\documentclass[11pt]{article}

\usepackage[T1]{fontenc}
\usepackage{lmodern}
\usepackage[a4paper,margin=30mm]{geometry}
\usepackage{amsmath,amssymb,amsthm,mathtools,mathrsfs}
\usepackage{microtype}
\usepackage{enumitem}
\usepackage{xcolor}
\usepackage[colorlinks=true,linkcolor=blue!45!black,citecolor=blue!45!black,
  urlcolor=blue!55!black]{hyperref}

\newtheorem{theorem}{Theorem}[section]
\newtheorem{proposition}[theorem]{Proposition}
\newtheorem{lemma}[theorem]{Lemma}
\newtheorem{corollary}[theorem]{Corollary}
\newtheorem{definition}[theorem]{Definition}
\newtheorem{remark}[theorem]{Remark}

\newcommand{\C}{\mathbb C}
\newcommand{\R}{\mathbb R}
\newcommand{\Q}{\mathbb Q}
\newcommand{\Z}{\mathbb Z}
\newcommand{\N}{\mathbb N}
\newcommand{\HH}{\mathbb H}
\newcommand{\cF}{\mathcal F}
\newcommand{\cG}{\mathcal G}
\newcommand{\cA}{\mathcal A}
\newcommand{\cK}{\mathcal K}
\newcommand{\cM}{\mathcal M}
\newcommand{\cR}{\mathcal R}
\newcommand{\cU}{\mathcal U}
\newcommand{\Fin}{\operatorname{Fin}}
\newcommand{\Fac}{\operatorname{Fac}}
\newcommand{\Span}{\operatorname{span}}
\newcommand{\dtrdeg}{\operatorname{dtrdeg}}
\newcommand{\card}[1]{\lvert #1\rvert}
\newcommand{\supp}{\operatorname{supp}}
\newcommand{\im}{\operatorname{Im}}
\newcommand{\D}{\partial}

\title{Universal Constraint--Orientation Product Geometry:\\
A Digamma Source and the Prime--Dirichlet Tail Law}
\author{Michael Haufschild \and GPT 5.6 Sol \and Fable 5 \and Grok 4.5}
\date{}

\begin{document}

\maketitle

\begin{abstract}
We construct a universal product between finite nonlinear differential constraint geometry and
analytic orientation.  In a Herglotz realization, a dense residual phase of positive prime-log
Cauchy sources simultaneously supports every finite localized Stanley--Reisner differential
matroid.  Every finite residue cube lies in one embedded differential field, while the Herglotz
condition selects exactly one vertex.  The finite products have a canonical inverse limit, the
product of the space of countable simplicial complexes with the prime-labelled Cantor cube, and
the all-positive sequence is its unique Herglotz point.

We then remove the residual source choice inside the Herglotz model itself.  The closed-form source
\[
  G_\star(z)=\psi\!\left(z+\tfrac12\right)+\gamma
\]
maps the upper half-plane strictly into itself, and the explicitly coded translates
$\psi(z+\alpha_q)+\gamma$, with $\alpha_q=1/(a(q)+2)$, form a jointly differentially
algebraically independent universal palette.  Finite residue reversals are rational gauges, and
the normally convergent signed Gauss series realizes every infinite orientation.  The all-positive
sequence is again exactly the Herglotz point, so the explicit source is selected intrinsically by
the geometry rather than chosen from a residual set.  More strongly, it is the unique meromorphic
solution of its half-integer unit-shift cocycle with the standard logarithmic normalization at
positive infinity.

At the analytic-completion level, every two absolutely summable residue carriers on the common
prime-log divisor generate the same compact-open closed rational Cauchy algebra.  This includes
all sources in the residual phase and every finite or infinite orientation.  Completion therefore
has an exact positive law: it forgets the source and orientation coordinates while retaining the
common atoms.

An explicit arithmetic model removes the remaining infinite-field ambiguity.  For a sign sequence
$\varepsilon$ on the primes, let
$P_\varepsilon(s)=\sum_p\varepsilon(p)p^{-s}$.  Over the finite-prime Dirichlet field, a finite
family $(P_{\varepsilon_j})$ is differentially algebraically independent exactly when the sign
classes $[\varepsilon_j]$ are linearly independent modulo finite support.  Consequently every
signed prime-zeta carrier is hypertranscendental, and two generated raw differential fields are
equal exactly when the signs agree eventually or are opposite eventually; otherwise their
compositum has differential transcendence degree two.  The same equivalence exactly classifies the
fields generated by their complete explicit integer-shift palettes.  Those palettes realize every
finite Stanley--Reisner constraint geometry for every orientation, while vertical positive
definiteness selects exactly the all-positive prime-zeta function.  Thus the inverse-limit field
coordinate is not left unspecified: it is the projective tail quotient of the orientation cube.
Every field fibre is countable and dense, the quotient topology is indiscrete, and every
nonconstant class-invariant map to a separated space is nowhere continuous.  The field relation is
continuously bireducible with the classical eventual-equality relation $E_0$, so the canonical
replacement for a representative selector is the pointed projective-tail groupoid.
The raw and completed statements form a two-scale classification, rather than an unresolved
same-field question.
\end{abstract}

\section{Introduction}

Hypertranscendence begins with the absence of an algebraic differential equation.  Since
H\"older's theorem for the gamma function \cite{Holder1886}, this absence has usually been treated
as a property of an individual special function.  The geometric question is different: what
structures can coexist on one field generated by a maximally free hypertranscendental source, and
which structures remain invisible after the source is forgotten?

Three established theories touch separate parts of this question.  Translation universality and
hypercyclicity provide functions with large orbits \cite{Birkhoff1929,Tsirivas2014}; hypercyclic
algebras organize nonlinear orbit phenomena \cite{BesEtAl2017,Bayart2018}; and
Stanley--Reisner theory records finite hereditary constraints by squarefree monomial ideals
\cite{Stanley1975}.  Herglotz--Nevanlinna theory, on the other hand, encodes positivity through
Cauchy transforms and their kernels \cite{RovnyakSakhnovich2004,Nedic2020}.  Rational changes of
Nevanlinna functions have an extensive theory of their own \cite{HassiWietsma2013}.  None of these
parent structures by itself asks whether arbitrary finite constraint geometry and analytic
orientation can vary independently inside one fixed differential field.

The first construction below answers that question without selecting a privileged point inside
its canonical generic phase.  For
every bounded gauge $u$ consider the positive prime-log Cauchy transform
\begin{equation}\label{eq:source-intro}
  g_u(z)=\sum_{p\ \mathrm{prime}}
    \frac{2^{-p}e^{u(p)}}{\log p-z}.
\end{equation}
The gauges for which one fixed collision-free palette of translates is differentially
algebraically independent form an intrinsically defined dense residual set $\cG_{\cU}$.  Every
$u\in\cG_{\cU}$ carries the full construction.  The palette contains the complete
facet--Laurent orbit required by every finite simplicial complex.  Facet restriction then realizes
the exact linear matroid of every localized Stanley--Reisner module as the nonlinear differential
matroid of analytic outputs.

The orientation coordinate is attached intrinsically to the prime-log carrier.  Fix an arbitrary
$u\in\cG_{\cU}$, write $\lambda_p=\log p$ and $w_p(u)=2^{-p}e^{u(p)}>0$.  For finite prime sets
$J\subseteq S\subset\mathcal P$, set
\begin{equation}\label{eq:chamber-intro}
  g_{u;S,J}(z)=g_u(z)-2\sum_{p\in J}\frac{w_p(u)}{\lambda_p-z}.
\end{equation}
Thus the Cauchy mass at $\lambda_p$ is $-w_p(u)$ for $p\in J$ and $+w_p(u)$ otherwise; $S$ records the
finite set of coordinates allowed to vary.  The correction is rational, so every finite chamber
generates the same embedded differential field.
The same rational-gauge invariance preserves the complete translated differential matroid.  In
contrast, a negative atom dominates the regular background near its own pole.  Consequently
$g_{u;S,J}$ is Herglotz exactly when $J=\varnothing$.

The Herglotz product theorem therefore has the quantifier order
\[
  \forall\,u\in\cG_{\cU}\quad\forall\,\Delta\quad\forall\,S
  \quad\forall\,J\subseteq S,
\]
with the source phase and palette fixed canonically before any finite simplicial complex or
orientation cube is chosen.  For each $u$, its embedded differential field is fixed throughout
the corresponding fiber.  The resulting geometry is
the product
\[
  \{\text{finite localized Stanley--Reisner constraint matroids}\}
  \times
  \{\text{finite oriented cubes}\}.
\]
The Herglotz cone provides a distinguished pointing of the second factor.  The finite restriction
maps have a canonical inverse limit, and summability realizes its complete orientation factor
analytically.  Finite rational residue changes are dense in a stronger sense: after compact-open
completion, every source in the phase and every infinite orientation generate the same
common-atom Cauchy algebra.  This completed collapse will be kept distinct from raw differential-
field equality.

The residual phase is not the endpoint.  A closed-form Herglotz realization starts from
\[
  G_\star(z)=\psi\!\left(z+\tfrac12\right)+\gamma.
\]
Its imaginary part has the positive Gauss expansion
\[
  \im G_\star(\sigma+it)
   =\sum_{n\geq0}\frac{t}{(n+\sigma+\tfrac12)^2+t^2}>0
   \qquad(t>0).
\]
The integer code used below gives shifts $\alpha_q=1/(a(q)+2)$ that are distinct modulo $\Z$.
The classical diagonal-difference theorem for Gamma translates then makes the digamma palette
$(\psi(z+\alpha_q)+\gamma)_{q\in\cU}$ jointly differentially algebraically independent.  This
explicit palette carries the complete constraint--orientation construction, while a signed Gauss
series gives a canonical analytic representative of every point of the infinite orientation cube.
The normalization is rigid: $G_\star$ is characterized uniquely by
$G(z+1)-G(z)=(z+\tfrac12)^{-1}$ and
$G(x)-\log x\to\gamma$ as $x\to+\infty$.

A third realization starts from the closed formula
\[
  P_+(s)=\sum_p p^{-s}.
\]
More generally, all signed prime-zeta functions $P_\varepsilon$ are hypertranscendental.  Their
complete differential matroid is exactly the linear matroid of the sign sequences modulo finite
support.  This identifies the raw-field boundary of the inverse limit: finite changes and global
sign reversal are precisely the same-field directions.  The positive orientation is again
intrinsic, now as the unique positive-definite vertical trace.  The arithmetic realization thus
supplies both an explicit universal source and an exact field stratification of the infinite
orientation cube.

Section~\ref{sec:upper} introduces the upper-end differential field.  Sections~\ref{sec:source}
and \ref{sec:palette} construct the universal source phase and palette.  Section~\ref{sec:gauge}
proves rational-gauge invariance.  Section~\ref{sec:orientation} proves the exact Herglotz chamber
law.  Sections~\ref{sec:sr} and \ref{sec:product} establish the constraint realization and the
product theorem.  Section~\ref{sec:limit} identifies the canonical inverse limit and its analytic
Herglotz pointing, together with the source--orientation invariant compact-open completion.
Section~\ref{sec:dirichlet} proves the explicit prime--Dirichlet tail law.
Section~\ref{sec:digamma} constructs the closed-form digamma Herglotz realization, after which the
arithmetic positive-definite pointing is identified.  The final section isolates the consequence
and its precise scope.

\section{Upper-end germs and differential independence}\label{sec:upper}

For $R\in\R$, write
\[
  \HH_R=\{z\in\C:\im z>R\}.
\]
Two meromorphic functions on upper half-planes are identified if they agree on a still higher
half-plane.  The resulting field of upper-end meromorphic germs is denoted by $\cK$.  It contains
$\C(z)$, is closed under $\D=d/dz$, and admits every complex translation
\[
  T_hf(z)=f(z+h),\qquad h\in\C.
\]
Indeed, if $f$ is meromorphic on $\HH_R$, then $T_hf$ is meromorphic on
$\HH_{R-\im h}$.  Differentiation and translation commute.

Let $K$ be a differential field of characteristic zero and let $B\subseteq K$ be a differential
subfield.  A family $(f_i)_{i\in I}$ is \emph{differentially algebraically independent over $B$}
if the jet family
\[
  \{\D^rf_i:i\in I,\ r\geq 0\}
\]
is algebraically independent over $B$.  We use the standard differential-algebraic terminology
of \cite{Kolchin1973}.  A single member of such a family is hypertranscendental over $B$.

For a set $E\subseteq K$, let $B\langle E\rangle$ denote the smallest differential subfield of
$K$ containing $B\cup E$.  Equality of fields below always means equality as subfields of $\cK$,
not merely an abstract isomorphism.

\section{A universal positive prime-log source}\label{sec:source}

Let $\mathcal P$ be the set of primes and let
\[
  \mathcal B=\ell^\infty(\mathcal P,\R)
\]
with the uniform norm.  For $u\in\mathcal B$, define
\begin{equation}\label{eq:gu}
  g_u(z)=\sum_{p\in\mathcal P}\frac{w_p(u)}{\lambda_p-z},
  \qquad
  \lambda_p=\log p,
  \qquad
  w_p(u)=2^{-p}e^{u(p)}.
\end{equation}
The series and all its derivatives converge locally uniformly away from the pole set.  In
particular, $g_u$ is meromorphic and
\begin{equation}\label{eq:jet-formula}
  \D^rg_u(z)=r!\sum_{p\in\mathcal P}
      \frac{w_p(u)}{(\lambda_p-z)^{r+1}}.
\end{equation}
Since every $w_p(u)$ is positive,
\begin{equation}\label{eq:base-herglotz}
  \im g_u(z)=\sum_{p\in\mathcal P}
    w_p(u)\frac{\im z}{\lvert\lambda_p-z\rvert^2}\geq 0
  \qquad(z\in\HH_0).
\end{equation}

The source theorem requires a finite controllability statement.  Its real nature is important:
the masses in \eqref{eq:gu} remain positive throughout the perturbation.

\begin{lemma}[Finite real Cauchy-jet controllability]\label{lem:control}
Let $\zeta_1,\ldots,\zeta_m$ be distinct points of the upper half-plane and let
$r_1,\ldots,r_m$ be positive integers.  The vectors
\[
  v(t)=\left((t-\zeta_j)^{-k}:1\leq j\leq m,\ 1\leq k\leq r_j\right),
  \qquad t\in\{\log p:p\in\mathcal P\},
\]
span $\prod_j\C^{r_j}$ as a real vector space.  A finite subset of these vectors already spans.
\end{lemma}

\begin{proof}
Suppose a real linear functional annihilates every $v(\log p)$.  After grouping real and
imaginary parts, there is a rational function
\[
  R(t)=\sum_{j=1}^m\sum_{k=1}^{r_j}\frac{c_{j,k}}{(t-\zeta_j)^k}
\]
such that
\[
  R(\log p)+\overline{R(\log p)}=0
\]
for every prime $p$.  The rational function
\[
  Q(t)=R(t)+\overline{R(\overline t)}
\]
therefore has infinitely many real zeros and is identically zero.  Its poles in the upper
half-plane are precisely the $\zeta_j$, whereas the second summand has poles only at the reflected
points $\overline{\zeta_j}$.  The principal part at each $\zeta_j$ must vanish, so every
$c_{j,k}$ is zero.  Thus the annihilating functional is zero.  Finite dimensionality yields a
finite spanning subfamily.
\end{proof}

\begin{lemma}[Coefficient-uniform pattern escape]\label{lem:pattern}
Fix distinct shifts $h_1,\ldots,h_m\in\C$, finite jet bounds, distinct differential monomials
$M_1,\ldots,M_s$ in those jets, and a nonnegative coefficient-degree bound $q$.  There are upper
evaluation points $z_1,\ldots,z_N$, where $N=s(q+1)$, such that the defect set
\begin{equation}\label{eq:defect}
  \left\{u\in\mathcal B:
  \det\left[z_a^dM_r\bigl((\D^jg_u(z_a+h_i))_{i,j}\bigr)
  \right]_{a,(r,d)}=0\right\}
\end{equation}
is closed and nowhere dense in $\mathcal B$.  Here $1\leq r\leq s$ and $0\leq d\leq q$.
\end{lemma}

\begin{proof}
Choose the $z_a$ on separated upper vertical half-lines so that all points $z_a+h_i$ are distinct.
The polynomials
\[
  (z,x)\longmapsto z^dM_r(x)
\]
are linearly independent.  Successively choosing evaluation pairs $(z_a,x_a)$ therefore produces
an invertible square evaluation matrix.  Equivalently, the determinant in
\eqref{eq:defect}, viewed as a polynomial in the finite collection of target jet coordinates, is
not the zero polynomial.

Fix $u\in\mathcal B$ and a uniform neighbourhood of $u$.  Apply
Lemma~\ref{lem:control} simultaneously to the distinct nodes $z_a+h_i$ and to all required jet
orders.  By \eqref{eq:jet-formula}, variations of finitely many real masses $w_p$ span the full
real jet space.  Hence the actual jet packet can be moved through a neighbourhood of its current
value.  The complement of the zero set of a nonzero complex polynomial is dense, so a sufficiently
small finite mass perturbation makes the determinant nonzero.  The perturbation may be chosen
smaller than every affected positive mass.  It therefore preserves positivity and has the form
$w_p'=2^{-p}e^{u'(p)}$ for a gauge $u'$ in the prescribed uniform neighbourhood.  Thus the defect
set has empty interior.

On a bounded gauge ball, the series in \eqref{eq:jet-formula} have a common summable majorant at
the finitely many evaluation nodes.  Every matrix entry, and hence the determinant, depends
continuously on $u$.  The defect set is closed, so it is nowhere dense.
\end{proof}

The determinant in Lemma~\ref{lem:pattern} excludes all complex coefficient vectors with the
given monomial support and degree bound at once.  This is what permits a countable Baire argument
despite the uncountability of $\C$.

\begin{theorem}[Countable-palette source phase]\label{thm:source}
Let $\Gamma$ be countable and let $s\colon\Gamma\to\C$ be injective.  There is a dense residual set
of gauges $u\in\mathcal B$ such that
\begin{equation}\label{eq:palette-di}
  \bigl(g_u(z+s(\gamma))\bigr)_{\gamma\in\Gamma}
\end{equation}
is differentially algebraically independent over $\C(z)$ in $\cK$.  Every such $g_u$ satisfies
the Herglotz inequality \eqref{eq:base-herglotz}.
\end{theorem}

\begin{proof}
A finite pattern consists of a finite injective shift packet from $\Gamma$, finite jet bounds, a
finite set of monomials, and a polynomial coefficient-degree bound.  There are only countably many
patterns.  By Lemma~\ref{lem:pattern}, the gauges that have full evaluation rank for a fixed
pattern form an open dense set.  The complete metric space $\mathcal B$ is a Baire space, so the
intersection of these open dense sets is dense and nonempty.

Choose $u$ in this intersection.  If \eqref{eq:palette-di} satisfied a nonzero differential
relation over $\C(z)$, clear the finitely many rational denominators.  The resulting relation has
polynomial coefficients of some finite degree, finitely many jet variables, and finitely many
distinct monomials.  It is one of the enumerated patterns.  Evaluation at the associated points
gives a full-rank matrix times its coefficient vector equal to zero.  The coefficient vector must
vanish, contradicting nontriviality.  Equation \eqref{eq:base-herglotz} gives the last assertion.
\end{proof}

\section{One collision-free palette for all finite complexes}\label{sec:palette}

Choose real numbers $\omega_0,\omega_1,\ldots$ linearly independent over $\Q$.  Let
$\Z^{(\N)}$ denote the finitely supported integer sequences and set
\[
  \cU=\Fin(\N)\times\Z^{(\N)}.
\]
Choose an injective map $\tau\colon\Fin(\N)\to\N$ with $\tau(\varnothing)=0$, and define
\begin{equation}\label{eq:universal-shift}
  s(A,m)=\sum_{n\geq 0}m_n\omega_n+i\tau(A).
\end{equation}

\begin{lemma}[Universal palette]\label{lem:palette}
The set $\cU$ is countable and the map $s\colon\cU\to\C$ is injective.  Moreover, for every finite
vertex set $V$ and every finite family of facets $\cF\subseteq\Fin(V)$, the map
\begin{equation}\label{eq:facet-embed}
  \cF\times\Z^V\longrightarrow\cU,
  \qquad
  (A,m)\longmapsto (e(A),\widetilde m),
\end{equation}
is injective after choosing an injection $e\colon V\to\N$ and extending $m$ by zero outside
$e(V)$.
\end{lemma}

\begin{proof}
Both finite subsets of $\N$ and finitely supported integer sequences are countable.  If
$s(A,m)=s(B,n)$, imaginary parts give $\tau(A)=\tau(B)$, hence $A=B$.  Real parts then give
$\sum_j(m_j-n_j)\omega_j=0$, and rational linear independence gives $m=n$.  The same two
arguments prove injectivity of \eqref{eq:facet-embed}.
\end{proof}

Apply Theorem~\ref{thm:source} to this fixed palette and define its intrinsic source phase
\begin{equation}\label{eq:source-phase}
  \cG_{\cU}=\left\{u\in\mathcal B:
  \bigl(g_u(z+s(q))\bigr)_{q\in\cU}
  \text{ is differentially algebraically independent over }\C(z)\right\}.
\end{equation}
This set is dense residual.  No point of it is selected.  For an arbitrary $u\in\cG_{\cU}$,
define
\begin{equation}\label{eq:fixed-field}
  \cF_u=\C(z)\left\langle g_u(z+s(q)):q\in\cU\right\rangle\subseteq\cK.
\end{equation}
The complete palette in \eqref{eq:fixed-field} is differentially algebraically independent over
$\C(z)$, and $g_u$ is Herglotz on $\HH_0$.  All subsequent assertions are uniform in
$u\in\cG_{\cU}$.

\section{Rational gauges do not change the field}\label{sec:gauge}

The following elementary differential-algebraic fact is the mechanism separating the field
coordinate from the orientation coordinate.

\begin{proposition}[Rational-gauge invariance]\label{prop:gauge}
Let $(f_i)_{i\in I}\subseteq\cK$ be differentially algebraically independent over $\C(z)$ and let
$r_i\in\C(z)$ for every $i$.  Then $(f_i+r_i)_{i\in I}$ is differentially algebraically
independent over $\C(z)$ and
\begin{equation}\label{eq:field-equality-general}
  \C(z)\langle f_i+r_i:i\in I\rangle
  =\C(z)\langle f_i:i\in I\rangle
\end{equation}
inside $\cK$.
\end{proposition}

\begin{proof}
Every differential relation involves a finite jet packet.  On its polynomial ring, substitution
is the automorphism
\[
  X_{i,k}\longmapsto X_{i,k}+\D^kr_i,
\]
whose inverse subtracts the same rational jet.  A nonzero polynomial relation after substitution
would therefore pull back to a nonzero relation among the jets of the $f_i$.  This proves
differential independence.  Since every derivative of every $r_i$ belongs to $\C(z)$, each side
of \eqref{eq:field-equality-general} contains the generators of the other.
\end{proof}

\section{The oriented Herglotz cube}\label{sec:orientation}

Fix $u\in\cG_{\cU}$.  For finite prime sets $J\subseteq S\subset\mathcal P$, define $g_{u;S,J}$ by
\eqref{eq:chamber-intro}.  The active coordinate set $S$ specifies a cube and $J$ is its set of
reversed prime masses.  Define the orientation index
\begin{equation}\label{eq:index}
  \kappa(S,J)=\card J.
\end{equation}

\begin{theorem}[Exact Herglotz chamber law]\label{thm:herglotz}
For arbitrary finite $J\subseteq S\subset\mathcal P$, the following are equivalent:
\begin{enumerate}[label=(\roman*)]
  \item $\im g_{u;S,J}(z)\geq 0$ for every $z\in\HH_0$;
  \item $J=\varnothing$;
  \item $\kappa(S,J)=0$.
\end{enumerate}
In particular, in the cube $J\subseteq S$, the all-positive vertex $J=\varnothing$ is the unique
Herglotz chamber.  The cube has $2^{\card S}$ vertices, and $\kappa(S,J)=\card J$ is the graph
distance from its all-positive vertex.
\end{theorem}

\begin{proof}
If $J=\varnothing$, then $g_{u;S,J}=g_u$, which is Herglotz by
\eqref{eq:base-herglotz}.  For reference, if $a\in\R$ and $z\in\HH_0$, then
\begin{equation}\label{eq:atom-positive}
  \im\frac{1}{a-z}=\frac{\im z}{\lvert a-z\rvert^2}>0.
\end{equation}
This proves (ii)$\Rightarrow$(i).

Conversely, choose $p\in J$.  Delete the $p$th source atom and absorb every other finite reversal
into the regular part.  Then
\[
  g_{u;S,J}(z)=G_p(z)-\frac{w_p(u)}{\lambda_p-z}.
\]
The function $G_p$ is continuous at $\lambda_p$: after deleting the $p$th term, normal convergence
of the remaining prime-log series gives a holomorphic germ there, and the finitely many other
reversals have distinct poles.  Hence $\im G_p(\lambda_p+iy)$ is bounded for sufficiently small
$y>0$.
But
\[
  \im\left(-\frac{w_p(u)}{\lambda_p-(\lambda_p+iy)}\right)=-\frac{w_p(u)}{y},
\]
so $\im g_{u;S,J}(\lambda_p+iy)<0$ for all sufficiently small $y>0$.  Thus (i) implies (ii).
Equivalence with (iii), the cube count, and the distance formula are immediate from finite-set
combinatorics.
\end{proof}

For $q\in\cU$, put
\begin{equation}\label{eq:chamber-palette}
  G_{u;S,J,q}(z)=g_{u;S,J}(z+s(q)).
\end{equation}
The translated correction
\[
  r_{S,J,q}(z)=-2\sum_{p\in J}
    \frac{w_p(u)}{\lambda_p-(z+s(q))}
\]
belongs to $\C(z)$.  Proposition~\ref{prop:gauge} therefore gives the following exact statement.

\begin{corollary}[One embedded field in every chamber]\label{cor:chamber-field}
For every finite $J\subseteq S\subset\mathcal P$, the palette $(G_{u;S,J,q})_{q\in\cU}$ is differentially
algebraically independent over $\C(z)$ and
\begin{equation}\label{eq:chamber-field}
  \C(z)\langle G_{u;S,J,q}:q\in\cU\rangle=\cF_u.
\end{equation}
\end{corollary}

\section{Finite constraint geometry by facet sewing}\label{sec:sr}

Let $V$ be finite and let $\Delta$ be a simplicial complex on $V$, specified by its finite facet
family $\Fac(\Delta)$.  Put $S_0=\C[\D]$ and
\[
  P=S_0[x_v:v\in V],\qquad t_v=1+x_v.
\]
The Stanley--Reisner ideal is
\begin{equation}\label{eq:sr-ideal}
  I_\Delta=(x_N:N\text{ is a nonface of }\Delta),
  \qquad x_N=\prod_{v\in N}x_v.
\end{equation}
It is enough to use the minimal nonfaces as generators.  We consider the localization
\begin{equation}\label{eq:localized-ring}
  A_\Delta=S_0[x_v:v\in V,t_v^{-1}:v\in V]/I_\Delta.
\end{equation}

For a facet $A$, define the restriction
\[
  \rho_A(x_v)=
  \begin{cases}
    x_v,&v\in A,\\
    0,&v\notin A,
  \end{cases}
  \qquad
  \rho_A(t_v)=
  \begin{cases}
    t_v,&v\in A,\\
    1,&v\notin A.
  \end{cases}
\]

\begin{lemma}[Exact facet coordinates]\label{lem:facet}
The product restriction map
\begin{equation}\label{eq:facet-map}
  A_\Delta\longrightarrow
  \prod_{A\in\Fac(\Delta)}S_0[t_v^{\pm1}:v\in A],
  \qquad P\longmapsto(\rho_A(P))_A,
\end{equation}
is injective.
\end{lemma}

\begin{proof}
Before localization, the standard facet-prime decomposition is
\begin{equation}\label{eq:facet-decomp}
  I_\Delta=\bigcap_{A\in\Fac(\Delta)}(x_v:v\notin A).
\end{equation}
Indeed, a monomial lies outside the right-hand side exactly when its support is contained in some
facet, equivalently when its support is a face.  Since monomials form an $S_0$-basis, the statement
holds coefficient by coefficient.  Localization at the finitely generated monoid of the $t_v$
is flat and commutes with this finite intersection.  In each facet quotient, every $t_v$ outside
the facet becomes $1$, while the remaining $t_v$ are invertible.  Thus the kernel of
\eqref{eq:facet-map} is zero.
\end{proof}

Fix an injection $e\colon V\to\N$.  For a facet $A$ and an integral vector $m\in\Z^V$ supported on
$A$, let
\begin{equation}\label{eq:facet-source}
  Y^{u;S,J}_{A,m}(z)=G_{u;S,J,(e(A),\widetilde m)}(z),
\end{equation}
where $\widetilde m$ is extended by zero to $\N$.  By Lemma~\ref{lem:palette} and
Corollary~\ref{cor:chamber-field}, the complete family in \eqref{eq:facet-source} is
differentially algebraically independent over $\C(z)$, hence also over $\C$.

Let $d=\prod_{v\in V}t_v^{b_v}$ be a common translation denominator.  Given numerator classes
$N_1,\ldots,N_r\in P/I_\Delta$, write the finite Laurent expansion
\begin{equation}\label{eq:row-expansion}
  \frac{\rho_A(N_i)}{\rho_A(d)}
  =\sum_{m\in\Z^A}L_{i,A,m}(\D)t^m,
  \qquad L_{i,A,m}\in\C[\D].
\end{equation}
Define the sewn analytic output
\begin{equation}\label{eq:sewn-output}
  \Phi^{u;S,J}_{\Delta,d}(N_i)
  =\sum_{A\in\Fac(\Delta)}\sum_{m\in\Z^A}
      L_{i,A,m}(\D)Y^{u;S,J}_{A,m}.
\end{equation}
Only finitely many terms occur.

\begin{lemma}[Differential operator rows]\label{lem:operator}
Let $(X_q)_{q\in Q}$ be differentially algebraically independent over $\C$, and let
$L_1,\ldots,L_r\in\C[\D]^{(Q)}$ have finite support.  Then the outputs
\[
  Z_i=\sum_qL_{i,q}(\D)X_q
\]
are differentially algebraically independent over $\C$ if and only if the rows $L_i$ are linearly
independent over $\C[\D]$.
\end{lemma}

\begin{proof}
A row dependence gives a nonzero constant-coefficient differential-linear relation among the
$Z_i$.  Conversely, restrict to the finite set of columns occurring in the rows.  Since
$\C[\D]$ is a Euclidean domain, Smith reduction transforms the row matrix by invertible row and
column operations into a diagonal matrix with one nonzero entry for each independent row.
Invertible row operations preserve the generated differential field, while invertible column
operations are differential changes of coordinates in the free differential field generated by
the $X_q$.  It remains to consider $L_i(\D)X_i$ with $L_i\neq0$.  In every putative differential
polynomial relation, the highest jet contributed by each $L_i$ has a nonzero leading coefficient;
substitution therefore has a unique highest-jet term and cannot vanish.  The diagonal outputs are
differentially algebraically independent.
\end{proof}

\begin{theorem}[Exact localized Stanley--Reisner differential matroid]\label{thm:sr}
For every $u\in\cG_{\cU}$, every finite $J\subseteq S\subset\mathcal P$, every finite simplicial complex $\Delta$, every common
translation denominator $d$, and every finite numerator family $(N_i)_{i\in I}$,
\begin{equation}\label{eq:matroid-law}
  \begin{aligned}
  &\bigl(\Phi^{u;S,J}_{\Delta,d}(N_i)\bigr)_{i\in I}
    \text{ is differentially algebraically independent over }\C\\
  &\hspace{38mm}\Longleftrightarrow
    (N_i)_{i\in I}\text{ is linearly independent over }\C[\D].
  \end{aligned}
\end{equation}
in $P/I_\Delta$.
\end{theorem}

\begin{proof}
The coordinates in \eqref{eq:row-expansion} are exactly the facet restrictions of the numerator
classes followed by the common Laurent exponent shift induced by $d$.  Division by one fixed
Laurent monomial is an injective reindexing of each finite coefficient row.  By
Lemma~\ref{lem:facet}, the resulting rows are linearly independent over $\C[\D]$ exactly when the
classes $N_i$ are.  Their analytic inputs \eqref{eq:facet-source} are differentially algebraically
independent.  Lemma~\ref{lem:operator} now gives \eqref{eq:matroid-law}.
\end{proof}

Thus every finite localized Stanley--Reisner linear matroid is realized, without extra nonlinear
differential dependencies, by analytic outputs in every orientation chamber.

\section{The universal constraint--orientation product}\label{sec:product}

We now collect the preceding results.  The source is quantified through its canonical phase rather
than selected from it.

\begin{theorem}[Universal constraint--orientation product]\label{thm:main}
There exist a countable injective translation palette $s\colon\cU\to\C$ and a canonically defined
dense residual phase $\cG_{\cU}\subset\mathcal B$ with the following simultaneous properties.
For every $u\in\cG_{\cU}$, let $g_u$ and $\cF_u\subseteq\cK$ be given by
\eqref{eq:source-intro} and \eqref{eq:fixed-field}.

For every such $u$, every finite $S\subseteq\mathcal P$, every chamber $J\subseteq S$, every finite simplicial complex
$\Delta$, every common translation denominator $d$, and every finite family of numerator classes
$(N_i)$:
\begin{enumerate}[label=(\alph*)]
  \item the complete chamber palette $(g_{u;S,J}(z+s(q)))_{q\in\cU}$ is differentially
    algebraically independent over $\C(z)$;
  \item it generates exactly $\cF_u$ as an embedded differential field;
  \item the sewn outputs realize the exact equivalence \eqref{eq:matroid-law};
  \item $g_{u;S,J}$ is Herglotz on the upper half-plane if and only if $J=\varnothing$;
  \item the orientation index is $\kappa(S,J)=\card J$, the cubical distance from the unique
    Herglotz vertex.
\end{enumerate}
Consequently, every field $\cF_u$ in the residual phase carries the Cartesian product of all
finite localized Stanley--Reisner differential matroids with all finite orientation cubes, and
the two coordinates can be varied independently.
\end{theorem}

\begin{proof}
Choose the universal palette of Lemma~\ref{lem:palette}; Theorem~\ref{thm:source} makes
$\cG_{\cU}$ dense residual.  Let $u\in\cG_{\cU}$ be arbitrary and define $\cF_u$ by
\eqref{eq:fixed-field}.  For each $S,J$, the
chamber correction and all its translated derivatives belong to $\C(z)$.  Proposition
\ref{prop:gauge} and Corollary~\ref{cor:chamber-field} prove (a) and (b).  The universal palette
contains the facet--Laurent orbit of every finite complex, so Theorem~\ref{thm:sr} proves (c) in
every chamber.  Theorem~\ref{thm:herglotz} proves (d) and (e).  The construction of the sewn rows
does not use $S$ or $J$, and the construction of the orientation cube does not use $\Delta$, $d$,
or $(N_i)$.  Hence the two projections vary independently, giving the stated product.
\end{proof}

\begin{corollary}[Arbitrary prescribed orientation index]\label{cor:index}
For every $u\in\cG_{\cU}$ and every $k\geq0$, the field $\cF_u$ has a representative which realizes every finite
localized Stanley--Reisner differential matroid, has orientation index exactly $k$, and is
Herglotz if and only if $k=0$.
\end{corollary}

\begin{proof}
Choose any $k$ primes and take their set for $S=J$ in Theorem~\ref{thm:main}.
\end{proof}

\begin{corollary}[Pointing law]\label{cor:pointing}
Fix a finite complex $\Delta$ and a nonempty finite support $S$.  No invariant that factors only
through
\begin{enumerate}[label=(\roman*)]
  \item the embedded differential field generated by the complete translated palette, and
  \item the finite sewn differential matroid of $\Delta$
\end{enumerate}
can recover whether the distinguished representative is Herglotz.  Herglotz positivity is instead
the unique basepoint of the orientation cube.
\end{corollary}

\begin{proof}
All chambers $J\subseteq S$ have the same embedded field by Theorem~\ref{thm:main}(b) and the same
sewn matroid by Theorem~\ref{thm:main}(c).  The chamber $J=\varnothing$ is Herglotz, whereas every
nonempty chamber is not, by Theorem~\ref{thm:main}(d).  Any invariant factoring through (i) and
(ii) is constant on this cube and therefore cannot equal its Herglotz indicator.  The positive
vertex is unique, so it gives a distinguished pointing.
\end{proof}

\section{The canonical profinite completion}\label{sec:limit}

The finite products in Theorem~\ref{thm:main} have an intrinsic completion.  Let
$\operatorname{Simp}(S)$ denote the set of downward-closed families of subsets of a finite vertex
set $S$.  For finite $S\subset\N$ and finite $T\subset\mathcal P$, put
\begin{equation}\label{eq:finite-stage}
  X_{S,T}=\operatorname{Simp}(S)\times\{+1,-1\}^{T}.
\end{equation}
Inclusion in either coordinate gives a bonding map by restriction.  The vertex and orientation
sets are deliberately independent: the directed set of stages consists of pairs $(S,T)$.

Write $\operatorname{Simp}(\N)$ for all downward-closed families of finite subsets of $\N$.

\begin{theorem}[Canonical constraint--orientation inverse limit]\label{thm:limit}
There is a canonical homeomorphism
\begin{equation}\label{eq:inverse-limit}
  \varprojlim_{S,T}X_{S,T}
  \cong
  \operatorname{Simp}(\N)\times\{+1,-1\}^{\mathcal P}.
\end{equation}
The second factor has a distinguished point $\boldsymbol{+}=(+1,+1,\ldots)$.  For every
$u\in\cG_{\cU}$ and every sign sequence $\varepsilon\in\{+1,-1\}^{\mathcal P}$, define
\begin{equation}\label{eq:infinite-carrier}
  g_{u,\varepsilon}(z)=
  \sum_{p\in\mathcal P}\frac{\varepsilon(p)w_p(u)}{\lambda_p-z}.
\end{equation}
This series converges normally on compact subsets of the complement of the prime-log divisor,
and the following assertions hold.
\begin{enumerate}[label=(\alph*)]
  \item If $\varepsilon$ is negative on a finite set $J$ and positive elsewhere, then
  $g_{u,\varepsilon}=g_{u;S,J}$ for every finite $S\supseteq J$.
  \item For an arbitrary $\varepsilon$, put
  $J_T=\{p\in T:\varepsilon(p)=-1\}$ at a finite stage $T$.  Then the finite chambers
  $g_{u;T,J_T}$ converge to $g_{u,\varepsilon}$, independently of the cofinal exhaustion of
  $\mathcal P$.
  \item At $\lambda_p$, the ordinary meromorphic residue is $-\varepsilon(p)w_p(u)\neq0$.
  \item The carrier $g_{u,\varepsilon}$ is Herglotz on $\HH_0$ if and only if
  $\varepsilon=\boldsymbol{+}$.
\end{enumerate}
Thus the all-positive vertices of all finite cubes form a compatible section whose limit is the
unique Herglotz point of the complete orientation fiber.
\end{theorem}

\begin{proof}
A global pair restricts to a coherent family of finite stages.  Conversely, let a coherent family
be given.  A finite set $A\subset\N$ is declared to be a face precisely when it is a face at the
stage $S=A$.  If $B\subseteq A$ and $A$ is a face, downward closure at stage $A$, followed by
coherence between the stages $B$ and $A$, shows that $B$ is a face.  The global sign at a prime
$p$ is read from the singleton stage $T=\{p\}$.  Coherence proves that these two reconstructions
are inverse to restriction.  They are coordinatewise continuous, hence give
\eqref{eq:inverse-limit}.  Equivalently, $\operatorname{Simp}(\N)$ is the closed subspace of
$\{0,1\}^{\operatorname{Fin}(\N)}$ cut out by the implications
$A\subseteq B$ and $x_B=1\Rightarrow x_A=1$; the other factor is the prime-labelled Cantor cube.

Boundedness of $u$ gives
\[
  \sum_p w_p(u)\leq e^{\lVert u\rVert_\infty}\sum_p2^{-p}<\infty.
\]
On a compact set away from the prime-log divisor, all sufficiently large $\lambda_p$ have
$\lvert\lambda_p-z\rvert\geq1$, while the finitely many remaining denominators have a positive
minimum.  The Weierstrass test therefore proves normal convergence of
\eqref{eq:infinite-carrier}.  Absolute convergence is unconditional, so every cofinal net of
finite partial sums has the same limit.  Expanding the sign gives
\[
  g_{u,\varepsilon}(z)=g_u(z)-2
    \sum_{\varepsilon(p)=-1}\frac{w_p(u)}{\lambda_p-z},
\]
which proves (a), (b), and the residue formula (c).

If every sign is positive, \eqref{eq:base-herglotz} gives the Herglotz inequality.  If
$\varepsilon(p)=-1$, delete the $p$th atom.  Normal convergence on the complement of the other
poles leaves a function $G_p$ holomorphic at $\lambda_p$, and
\[
  g_{u,\varepsilon}(z)=G_p(z)-\frac{w_p(u)}{\lambda_p-z}.
\]
At $z=\lambda_p+iy$, the second term has imaginary part $-w_p(u)/y$, whereas
$\im G_p(\lambda_p+iy)$ remains bounded as $y\downarrow0$.  Hence
$g_{u,\varepsilon}$ fails the Herglotz inequality.  This proves (d) and the uniqueness statement.
\end{proof}

The inverse limit has a second, coarser analytic level on which every orientation belongs to one
canonical object.  For an absolutely summable real family $a=(a_p)_{p\in\mathcal P}$ put
\begin{equation}\label{eq:summable-carrier}
  C_a(z)=\sum_{p\in\mathcal P}\frac{a_p}{\lambda_p-z},\qquad z\in\HH_0.
\end{equation}
Let $r_p(z)=(\lambda_p-z)^{-1}$, let $\cA_0$ be the complex algebra generated by the
$r_p$, and define
\begin{equation}\label{eq:completed-ring}
  \widehat{\cA}_a=
  \overline{\operatorname{alg}(\cA_0,C_a)}^{\,\mathrm{co}}
  \subset \mathcal O(\HH_0),
\end{equation}
where the closure is taken in the compact-open topology.

\begin{theorem}[Summable-residue completion law]\label{thm:summable-gauge}
Let $a,b\in\ell^1(\mathcal P,\R)$.  Then
\begin{equation}\label{eq:finite-transition-limit}
  C_a+\sum_{p\in F}(b_p-a_p)r_p
  \longrightarrow C_b
\end{equation}
locally uniformly on $\HH_0$ as the finite set $F$ increases to $\mathcal P$.  Consequently
\begin{equation}\label{eq:completed-ring-equality}
  \widehat{\cA}_a=\widehat{\cA}_b.
\end{equation}
More generally, if $\cR\subset\mathcal O(\HH_0)$ is any compact-open closed subring containing all
complex constants and all $r_p$, then
\begin{equation}\label{eq:closed-ring-indistinguishability}
  C_a\in\cR\quad\Longleftrightarrow\quad C_b\in\cR.
\end{equation}
In particular, the minimal compact-open closed differential subring containing the common atoms
and one summable carrier is independent of that carrier.
\end{theorem}

\begin{proof}
The difference $b-a$ is absolutely summable.  On a compact set $K\subset\HH_0$, all sufficiently
large $\lambda_p$ satisfy $|\lambda_p-z|\geq1$ for $z\in K$, and the finitely many remaining
denominators have a positive minimum on $K$.  The Weierstrass test therefore gives locally uniform
and unconditional convergence of
\[
  \sum_p(b_p-a_p)r_p=C_b-C_a,
\]
which proves \eqref{eq:finite-transition-limit}.  Every finite approximant belongs to the algebra
generated by $\cA_0$ and $C_a$, so closedness gives $C_b\in\widehat{\cA}_a$.  Interchanging $a$
and $b$ yields the reverse containment and proves \eqref{eq:completed-ring-equality}.  The same
argument inside $\cR$ proves \eqref{eq:closed-ring-indistinguishability}.  A closed differential
subring is a special case; intersecting all such subrings gives the final assertion.
\end{proof}

\begin{corollary}[Source--orientation completion collapse]
\label{cor:source-orientation-completion}
For arbitrary $u,v\in\cG_{\cU}$ and arbitrary sign sequences $\varepsilon,\delta$, set
\[
  a_p=\varepsilon(p)2^{-p}e^{u(p)},\qquad
  b_p=\delta(p)2^{-p}e^{v(p)}.
\]
Then $a,b\in\ell^1(\mathcal P,\R)$ and
\[
  \widehat{\cA}_{u,\varepsilon}:=\widehat{\cA}_a
  =\widehat{\cA}_b=: \widehat{\cA}_{v,\delta}.
\]
Thus the compact-open completed rational Cauchy algebra is constant on the entire product
$\cG_{\cU}\times\{+1,-1\}^{\mathcal P}$, including every infinitely reversed sign sequence.
\end{corollary}

\begin{proof}
Boundedness of $u$ and $v$ reduces both absolute sums to constant multiples of $\sum_p2^{-p}$.
Apply Theorem~\ref{thm:summable-gauge}.
\end{proof}

\begin{remark}[Two-scale field law]\label{rem:field-boundary}
The completed equality is not a claim that all raw differential fields coincide.  The completion
forgets the residue tail because finite rational transitions are dense.  The raw arithmetic fields
retain exactly the projective tail class, as Corollaries~\ref{cor:raw-field-quotient} and
\ref{cor:complete-palette-field-quotient} below show.  Hence the infinite inverse limit has an
exact two-scale geometry: one completed analytic algebra, stratified before completion by
projective tail differential fields.
\end{remark}

\section{The explicit prime--Dirichlet tail law}\label{sec:dirichlet}

The Cauchy realization proves the geometry, but its source phase is residual.  We now give an
arithmetic realization with a closed-form source and identify the raw differential field at every
point of the infinite orientation cube.

Let $\mathscr A$ be the complex algebra of arithmetic functions under Dirichlet convolution.  If
$a\in\mathscr A$, let $[\supp a]$ be the set of primes dividing some integer in its support, and
let $\mathscr S\subset\mathscr A$ be the subalgebra of functions for which $[\supp a]$ is finite.
Write $\mathscr E_{\mathrm{fin}}$ for the differential field of right-half-plane meromorphic germs
generated by $\C(s)$ and the convergent Dirichlet series attached to $\mathscr S$.  Thus each
individual element of $\mathscr E_{\mathrm{fin}}$ sees only finitely many primes, although no
uniform finite prime set is fixed for the whole field.

For $\varepsilon\in\{+1,-1\}^{\mathcal P}$ define, on $\Re s>1$,
\begin{equation}\label{eq:signed-prime-zeta}
  P_\varepsilon(s)=\sum_{p\in\mathcal P}\varepsilon(p)p^{-s}.
\end{equation}
Let
\[
  V_{\mathrm{tail}}=
  \C^{\mathcal P}/\C^{(\mathcal P)},
\]
where $\C^{(\mathcal P)}$ is the space of finitely supported sequences.  We write
$[\varepsilon]$ for the image of a sign sequence in $V_{\mathrm{tail}}$.

\begin{theorem}[Shifted tail differential matroid]\label{thm:dirichlet-tail}
Let $A\subset\N$ be finite and nonempty.  For every finite family
$\varepsilon_1,\ldots,\varepsilon_n$ of prime sign sequences, the following are equivalent:
\begin{enumerate}[label=(\roman*)]
  \item the shifted carriers
    $\bigl(P_{\varepsilon_j}(s+a)\bigr)_{1\leq j\leq n,\ a\in A}$ are differentially
    algebraically independent over $\mathscr E_{\mathrm{fin}}$;
  \item the classes $[\varepsilon_1],\ldots,[\varepsilon_n]$ are linearly independent in
    $V_{\mathrm{tail}}$.
\end{enumerate}
Consequently every $P_\varepsilon$ is hypertranscendental over
$\mathscr E_{\mathrm{fin}}$.
\end{theorem}

\begin{proof}
Let $a_\varepsilon\in\mathscr A$ be supported on the primes, with
$a_\varepsilon(p)=\varepsilon(p)$.  The logarithmic derivation
$\partial_La(n)=(\log n)a(n)$ corresponds, up to sign, to differentiation of the Dirichlet
series, while multiplication of coefficients by $n^{-a}$ corresponds to translation by $a$.
Hence the coefficient at a prime $p$ of the $k$th logarithmic jet of the $a$-shift is
\begin{equation}\label{eq:prime-jet-column}
  \varepsilon_j(p)p^{-a}(\log p)^k.
\end{equation}

Assume first that the tail classes are independent, and fix a jet order $m$.  A sign pattern
$b\in\{+1,-1\}^n$ is called recurrent when
\[
  \{p:(\varepsilon_1(p),\ldots,\varepsilon_n(p))=b\}
\]
is infinite.  Tail independence is equivalent to the recurrent patterns spanning $\C^n$: an
annihilating linear form can be nonzero only on the union of the nonrecurrent pattern fibres, and
that union is finite.

Suppose a linear combination of the rows in \eqref{eq:prime-jet-column}, for $a\in A$ and
$0\leq k\leq m$, vanishes outside finitely many primes.  On an infinite fibre with pattern $b$,
it has the form
\begin{equation}\label{eq:shifted-exp-polynomial}
  \sum_{a\in A}p^{-a}Q_{a,b}(\log p)=0,
\end{equation}
where each $Q_{a,b}$ is a polynomial.  If some $Q_{a,b}$ is nonzero, choose the least corresponding
$a_0$.  Multiplication by $p^{a_0}$ gives
\[
  Q_{a_0,b}(\log p)+
  \sum_{a>a_0}p^{-(a-a_0)}Q_{a,b}(\log p)=0.
\]
Along the infinite prime fibre, $p\to\infty$ and the second term tends to zero.  A nonzero
polynomial at $\log p\to\infty$ cannot tend to zero, a contradiction.  Thus every $Q_{a,b}$
vanishes.  Its coefficient at every power $(\log p)^k$ annihilates every recurrent pattern, and
the spanning property makes every original coefficient zero.  Hence the complete finite
shift--jet rectangle remains linearly independent after deletion of an arbitrary finite prime set.

It follows that, outside any prescribed finite set of primes, one can choose
$n\card{A}(m+1)$ primes for which the square matrix \eqref{eq:prime-jet-column} is invertible.  For a prime
$q$, the basic derivation of the arithmetic-function algebra is
\[
  (\partial_q a)(r)=v_q(rq)a(rq).
\]
The value at $1$ of the Jacobian determinant of the selected jet functions under these basic
derivations is precisely the preceding nonzero determinant.  The Shapiro--Sparer Jacobian
criterion \cite{ShapiroSparer1986}, in the strengthened form recorded as Theorem~3.5 by Pong
\cite{Pong2016}, makes the jet rectangle algebraically independent over the common kernel of the
selected derivations.  Given a proposed relation over $\mathscr S$, choose the primes after
deleting the finite union of prime supports of its coefficients.  All coefficients then lie in
that common kernel, so no relation exists.  Clearing denominators gives the same conclusion over
the fraction field.

The formal Dirichlet-series isomorphism sends $\partial_L$ to $-d/ds$.  For convergent series,
Dirichlet-series uniqueness transfers the result to \eqref{eq:signed-prime-zeta}.  To pass from the
fraction field of $\mathscr S$ to $\mathscr E_{\mathrm{fin}}$, clear all field denominators in a
putative relation and then the common rational denominator in $s$.  Its coefficients become finite
sums of terms $s^rB(s)$ with $B$ a convergent Dirichlet series from $\mathscr S$.  The standard
polynomial-weighted uniqueness argument along a real ray separates the powers of $s$ and would
produce a nonzero relation over $\mathscr S$, which has just been excluded.  This is the classical
formal-to-analytic passage used in \cite{ShapiroSparer1986,Pong2016}.  Since $m$ was arbitrary,
(ii) implies (i).

Conversely, if $\sum_jc_j\varepsilon_j$ has finite support for a nonzero vector $c$, choose any
$a\in A$.  Then
\[
  \sum_jc_jP_{\varepsilon_j}(s+a)
    =\sum_{p\in F}\left(\sum_jc_j\varepsilon_j(p)\right)p^{-a}p^{-s}
    \in\mathscr E_{\mathrm{fin}}
\]
for a finite prime set $F$.  This is a nontrivial differential algebraic relation, proving the
reverse implication.  A single sign sequence has infinite support modulo finite sequences, which
gives the final assertion.
\end{proof}

\begin{corollary}[Exact raw-field quotient]\label{cor:raw-field-quotient}
Put
\[
  \mathscr K_\varepsilon=
  \mathscr E_{\mathrm{fin}}\langle P_\varepsilon\rangle.
\]
For two orientations $\varepsilon,\delta$,
\begin{equation}\label{eq:raw-field-classification}
  \mathscr K_\varepsilon=\mathscr K_\delta
  \quad\Longleftrightarrow\quad
  \begin{cases}
  \varepsilon(p)=\delta(p)&\text{for all but finitely many }p,\quad\text{or}\cr
  \varepsilon(p)=-\delta(p)&\text{for all but finitely many }p.
  \end{cases}
\end{equation}
If neither alternative holds, $P_\varepsilon$ and $P_\delta$ are differentially algebraically
independent over $\mathscr E_{\mathrm{fin}}$, and their compositum has differential transcendence
degree two.
\end{corollary}

\begin{proof}
Two sign classes in $V_{\mathrm{tail}}$ are linearly dependent exactly when the signs are eventually
equal or eventually opposite.  Indeed, if
$a\varepsilon(p)+b\delta(p)=0$ off a finite set and both the agreement and disagreement sets were
infinite, evaluating on the two sets would give $a+b=0$ and $a-b=0$, contrary to
$(a,b)\neq(0,0)$.  The converse follows from the relations
$\varepsilon-\delta$ and $\varepsilon+\delta$.

In either eventual case, $P_\varepsilon\mp P_\delta$ is a finite Dirichlet polynomial in
$\mathscr E_{\mathrm{fin}}$, so the generated differential fields are equal.  Outside those cases,
Theorem~\ref{thm:dirichlet-tail} gives differential independence.  If the two generated fields
were equal, one carrier would be a differential rational function of finitely many jets of the
other, producing a forbidden algebraic relation.  The differential-transcendence-degree statement
is immediate.
\end{proof}

The source and its universal constraint palette are now completely explicit.  Enumerate the
primes increasingly as $q_0,q_1,\ldots$.  For a finite $A\subset\N$ and
$m\in\Z^{(\N)}$, put
\[
  c(A)=\sum_{j\in A}2^j,
  \qquad
  c(m)=\prod_{m_j>0}q_{2j}^{m_j}
       \prod_{m_j<0}q_{2j+1}^{-m_j},
\]
and let $\pi(x,y)=\frac12(x+y)(x+y+1)+y$ be the Cantor pairing function.  Then
\begin{equation}\label{eq:explicit-shift-code}
  a(A,m)=\pi(c(A),c(m)-1)
\end{equation}
is an injective map $\cU\to\N$ and sends $(\varnothing,0)$ to $0$.

\begin{theorem}[Explicit universal arithmetic source]\label{thm:explicit-source}
For every orientation $\varepsilon$, the countable family
\begin{equation}\label{eq:explicit-palette}
  \bigl(P_\varepsilon(s+a(q))\bigr)_{q\in\cU}
\end{equation}
is differentially algebraically independent over $\mathscr E_{\mathrm{fin}}$.  Replacing the
Cauchy palette by \eqref{eq:explicit-palette} in the facet-sewing construction realizes every
finite localized Stanley--Reisner differential matroid.  If two orientations differ at finitely
many primes, their complete palettes differ coordinatewise by elements of
$\mathscr E_{\mathrm{fin}}$ and therefore generate the same embedded differential field.
\end{theorem}

\begin{proof}
Unique factorization and injectivity of Cantor pairing prove that \eqref{eq:explicit-shift-code} is
injective.  Any finite subfamily of \eqref{eq:explicit-palette} has a finite nonempty set of distinct
shifts.  Theorem~\ref{thm:dirichlet-tail}, applied to the one-element orientation family, makes that
subfamily differentially algebraically independent.  Differential algebraic dependence is
finitary, so the complete countable palette is independent.  The proofs of Lemmas~\ref{lem:operator} and
\ref{lem:facet} and Theorem~\ref{thm:sr} use only differential independence of the source packet,
so they apply without change.  Finally, a finite sign change gives, for every $q$,
\[
  P_\varepsilon(s+a(q))-P_\delta(s+a(q))
   =\sum_{p\in F}(\varepsilon(p)-\delta(p))p^{-a(q)}p^{-s},
\]
which belongs to $\mathscr E_{\mathrm{fin}}$.  Differential base translation preserves both the
field and the sewn matroid.
\end{proof}

The preceding finite-shift theorem also determines the field generated by the entire countable
palette, because every differential-rational expression involves only finitely many palette
coordinates.  Put
\begin{equation}\label{eq:complete-arithmetic-field}
  \mathscr L_\varepsilon=
  \mathscr E_{\mathrm{fin}}
  \left\langle P_\varepsilon(s+a(q)):q\in\cU\right\rangle.
\end{equation}

\begin{corollary}[Exact complete-palette field quotient]
\label{cor:complete-palette-field-quotient}
For two orientations $\varepsilon,\delta$,
\begin{equation}\label{eq:complete-palette-field-classification}
  \mathscr L_\varepsilon=\mathscr L_\delta
  \quad\Longleftrightarrow\quad
  \begin{cases}
  \varepsilon(p)=\delta(p)&\text{for all but finitely many }p,\quad\text{or}\cr
  \varepsilon(p)=-\delta(p)&\text{for all but finitely many }p.
  \end{cases}
\end{equation}
If neither alternative holds, every finite subfamily of the combined palettes of $\varepsilon$
and $\delta$ is differentially algebraically independent over $\mathscr E_{\mathrm{fin}}$.
\end{corollary}

\begin{proof}
In either eventual case, for every $q\in\cU$ one has
\[
  P_\varepsilon(s+a(q))\mp P_\delta(s+a(q))
  =\sum_{p\in F}(\varepsilon(p)\mp\delta(p))p^{-a(q)}p^{-s}
  \in\mathscr E_{\mathrm{fin}}
\]
for one finite prime set $F$.  The two complete palettes therefore generate the same field.

Suppose neither eventual relation holds.  Corollary~\ref{cor:raw-field-quotient} says that the two
tail classes are linearly independent.  Any finite subfamily of the combined palettes is contained
in
\[
  \bigl(P_\varepsilon(s+a),P_\delta(s+a)\bigr)_{a\in A}
\]
for a finite nonempty set $A$ of coded shifts.  Theorem~\ref{thm:dirichlet-tail} makes this finite
rectangle differentially algebraically independent, proving the last assertion.  If
$\mathscr L_\varepsilon=\mathscr L_\delta$, one coordinate of the $\varepsilon$-palette would be a
differential-rational expression in finitely many coordinates of the $\delta$-palette.  Clearing
its denominator would contradict precisely that finite independence.  Thus the fields are
unequal.
\end{proof}

\begin{corollary}[Finite--cofinite closure of the canonical source]
\label{cor:finite-cofinite-source-field}
Let
\[
  N_\varepsilon=\{p\in\mathcal P:\varepsilon(p)=-1\}.
\]
Then the complete palette of an arbitrary orientation belongs to the same embedded differential
field as the canonical all-positive source precisely in the finite--cofinite cases:
\begin{equation}\label{eq:finite-cofinite-source-field}
  \mathscr L_\varepsilon=\mathscr L_{\boldsymbol{+}}
  \quad\Longleftrightarrow\quad
  N_\varepsilon\text{ is finite or }\mathcal P\setminus N_\varepsilon\text{ is finite}.
\end{equation}
If both sets are infinite, every finite subfamily of the two complete palettes is differentially
algebraically independent over $\mathscr E_{\mathrm{fin}}$.  Thus an infinite reversal has exactly
two possible field behaviours: a cofinite reversal is the global-sign direction modulo a finite
gauge and preserves the field, whereas a genuinely mixed infinite tail creates a transverse
differential-field direction.
\end{corollary}

\begin{proof}
Apply Corollary~\ref{cor:complete-palette-field-quotient} with
$\delta=\boldsymbol{+}$.  Eventual equality with $\boldsymbol{+}$ says exactly that
$N_\varepsilon$ is finite, and eventual opposition says exactly that its complement is finite.
The joint-independence assertion is the remaining case of that corollary.
\end{proof}

\begin{theorem}[Canonical projective-tail field bundle]
\label{thm:projective-tail-field-bundle}
On the orientation cube $\Omega=\{\pm1\}^{\mathcal P}$ define
\begin{equation}\label{eq:projective-tail-relation}
  \varepsilon\sim_{\mathrm{pt}}\delta
  \quad\Longleftrightarrow\quad
  \{p:\varepsilon(p)\ne\delta(p)\}\text{ is finite}
  \quad\text{or}\quad
  \{p:\varepsilon(p)=\delta(p)\}\text{ is finite}.
\end{equation}
Then $\sim_{\mathrm{pt}}$ is an equivalence relation, canonically equal to proportionality of the
nonzero tail lines $\C[\varepsilon]\subset V_{\mathrm{tail}}$.  The assignment
\begin{equation}\label{eq:projective-tail-field-map}
  \Phi:\Omega/\!\sim_{\mathrm{pt}}\ \longrightarrow
  \{\mathscr L_\varepsilon:\varepsilon\in\Omega\},
  \qquad [\varepsilon]_{\mathrm{pt}}\longmapsto\mathscr L_\varepsilon,
\end{equation}
is a bijection.  Every finite orientation chamber, extended by positive signs outside its finite
stage, belongs to the distinguished fibre $[\boldsymbol{+}]_{\mathrm{pt}}$; the same fibre also
contains the cofinite reversals.  A mixed infinite tail belongs to a different fibre, and its
palette is jointly differentially algebraically independent from the positive palette on every
finite subfamily.

Thus the field carried by the inverse limit is canonical as a bundle, not as one raw field.  Its
base is the projective-tail quotient and its distinguished finite-chamber fibre is
$\mathscr L_{\boldsymbol{+}}$.  Under the common orientation indexing, the compact-open Cauchy
realization sends every fibre to the single algebra of
Corollary~\ref{cor:source-orientation-completion}.  The point
$\boldsymbol{+}$ is nevertheless selected uniquely inside the full orientation cube by Herglotz
positivity in the digamma realization and by positive definiteness in the arithmetic realization.
\end{theorem}

\begin{proof}
The projective characterization in the proof of Corollary~\ref{cor:raw-field-quotient} shows that
\eqref{eq:projective-tail-relation} is precisely equality of the lines spanned by the nonzero
classes $[\varepsilon]$ and $[\delta]$ in $V_{\mathrm{tail}}$; hence it is an equivalence relation.
Corollary~\ref{cor:complete-palette-field-quotient} says exactly that $\Phi$ is well defined and
injective, while surjectivity holds by its displayed codomain.  Finite and cofinite membership in
the distinguished fibre and finite-subfamily independence are
Corollary~\ref{cor:finite-cofinite-source-field}.  Completion collapse is
Corollary~\ref{cor:source-orientation-completion}; the two uniqueness assertions are
Theorems~\ref{thm:digamma-source} and \ref{thm:positive-definite-pointing}.
\end{proof}

\begin{theorem}[Topological opacity of the projective-tail field bundle]
\label{thm:projective-tail-opacity}
Equip $\Omega=\{\pm1\}^{\mathcal P}$ with its product topology.  Then:
\begin{enumerate}
  \item every fibre of the field map $\varepsilon\mapsto\mathscr L_\varepsilon$ is countable and
    dense in $\Omega$;
  \item the quotient topology on $\Omega/\!\sim_{\mathrm{pt}}$ is indiscrete;
  \item if $Y$ is a $T_1$ space and $F:\Omega\to Y$ is constant on the field fibres, then
    continuity of $F$ at one orientation forces $F$ to be constant everywhere.  In particular,
    every nonconstant field-fibre invariant is nowhere continuous.
\end{enumerate}
After an enumeration of the primes, $\sim_{\mathrm{pt}}$ is continuously bireducible, and hence
Borel bireducible, with the classical eventual-equality relation $E_0$ on $2^{\N}$.  Consequently
it is nonsmooth: there
is no Borel complete invariant with values in a standard Borel space and no Borel transversal
selecting one orientation from every embedded-field fibre.
\end{theorem}

\begin{proof}
The fibre through $\varepsilon$ consists of the finite modifications of $\varepsilon$ and of
$-\varepsilon$.  There are only countably many such modifications.  Every nonempty basic cylinder
prescribes only finitely many signs, so it contains a finite modification of $\varepsilon$; hence
every fibre is dense.

If a nonempty open saturated set contains one basic cylinder, every field fibre meets that
cylinder.  Saturation then contains every point of $\Omega$, proving that the quotient topology is
indiscrete.  For the pointwise assertion, suppose that $F$ is continuous at $\varepsilon_0$.  If
$U$ is a neighbourhood of $F(\varepsilon_0)$, then $F^{-1}(U)$ contains a neighbourhood of
$\varepsilon_0$.  The dense fibre through any $\delta$ meets this neighbourhood.  Fibre invariance
therefore gives $F(\delta)\in U$.  Thus $F(\delta)$ belongs to every neighbourhood of
$F(\varepsilon_0)$; the $T_1$ property forces $F(\delta)=F(\varepsilon_0)$.

Identify orientations with Boolean sequences.  Define
\[
  D(x)(n)=x(n)\mathbin{\mathtt{xor}}x(n+1).
\]
If $x$ and $y$ are eventually equal or eventually complementary, then $D(x)$ and $D(y)$ are
eventually equal.  Conversely, eventual equality of $D(x)$ and $D(y)$ says that
$x(n)\mathbin{\mathtt{xor}}y(n)$ is eventually constant, giving eventual equality or eventual
complementarity.  Hence $D$ is a continuous reduction of $\sim_{\mathrm{pt}}$ to $E_0$.
For the reverse reduction put
\[
  A(x)(2n)=x(n),\qquad A(x)(2n+1)=0.
\]
The anchored odd coordinates exclude eventual complementarity, so
$A(x)\sim_{\mathrm{pt}}A(y)$ exactly when $xE_0y$.  The maps $D$ and $A$ are coordinatewise
continuous.  The final assertions now follow from the classical nonsmoothness theory of $E_0$;
see Dougherty, Jackson and Kechris~\cite{DoughertyJacksonKechris1994}.
\end{proof}

\section{The explicit digamma Herglotz realization}\label{sec:digamma}

The arithmetic source makes the field coordinate explicit, but its positive pointing is expressed
by positive definiteness rather than the Herglotz cone.  We now give a closed-form source inside
the Herglotz realization itself.  Let $\psi=\Gamma'/\Gamma$ be the digamma function, let $\gamma$
be Euler's constant, and put
\begin{equation}\label{eq:digamma-source}
  G_\star(z)=\psi\!\left(z+\frac12\right)+\gamma.
\end{equation}
Retain the injective code $a:\cU\to\N$ from \eqref{eq:explicit-shift-code}, and define
\begin{equation}\label{eq:digamma-palette-code}
  \alpha_q=\frac1{a(q)+2},
  \qquad
  \tau_q=\alpha_q-\frac12
  \qquad(q\in\cU).
\end{equation}
The code sends $(\varnothing,0)$ to zero, so the unshifted source occurs in its own palette.

Enumerate the primes increasingly as $q_0,q_1,\ldots$ and write $\nu(q_n)=n$.  For a finite set
$J$ of primes, define the normalized residue reversal
\begin{equation}\label{eq:finite-digamma-orientation}
  G_{\star,J}(z)
  =G_\star(z)-2\sum_{p\in J}
    \left(\frac1{\nu(p)+1}-\frac1{\nu(p)+z+\tfrac12}\right).
\end{equation}
The real constants in this normalization are immaterial to both the differential field and the
Herglotz condition; they make finite truncations agree exactly with the infinite signed series
below.

\begin{theorem}[Closed-form Herglotz universal source]\label{thm:digamma-source}
The source \eqref{eq:digamma-source} has the following properties.
\begin{enumerate}[label=\textup{(\roman*)}]
  \item The countable palette
  \begin{equation}\label{eq:digamma-palette}
    \bigl(G_\star(z+\tau_q)\bigr)_{q\in\cU}
      =\bigl(\psi(z+\alpha_q)+\gamma\bigr)_{q\in\cU}
  \end{equation}
  is differentially algebraically independent over $\C(z)$.
  \item The function $G_\star$ maps the upper half-plane strictly into itself.
  \item For every finite $J$, the palette
  $(G_{\star,J}(z+\tau_q))_{q\in\cU}$ is differentially algebraically independent and generates
  the same embedded differential field as \eqref{eq:digamma-palette}.  Moreover,
  $G_{\star,J}$ is Herglotz if and only if $J=\varnothing$.
  \item Every orientation $\varepsilon\in\{\pm1\}^{\mathcal P}$ has the canonical representative
  \begin{equation}\label{eq:signed-digamma-completion}
    G_\varepsilon(z)=\sum_{n=0}^{\infty}\varepsilon(q_n)
      \left(\frac1{n+1}-\frac1{n+z+\tfrac12}\right).
  \end{equation}
  This series converges locally normally away from the half-integer pole set, has residue
  $-\varepsilon(q_n)$ at $-n-\tfrac12$, and is Herglotz exactly when
  $\varepsilon=\boldsymbol{+}$.  If $\varepsilon^{(N)}$ agrees with $\varepsilon$ on
  $q_0,\ldots,q_{N-1}$ and is positive thereafter, then
  $G_{\varepsilon^{(N)}}\to G_\varepsilon$ locally normally away from the poles.
\end{enumerate}
Consequently the single explicit source $G_\star$ supports every finite localized
Stanley--Reisner differential matroid in every finite orientation chamber, and
\eqref{eq:signed-digamma-completion} realizes the full orientation inverse limit with a unique
Herglotz pointing.
\end{theorem}

\begin{proof}
The numbers $\alpha_q$ lie in $(0,1/2]$.  Injectivity of $a$ makes them distinct, and two distinct
members cannot differ by an integer.  Hardouin's diagonal-difference theorem
\cite[Theorem~4.2 and Corollary~4.4]{Hardouin2008} therefore applies to every finite family of
Gamma translates $\Gamma(z+\alpha_q)$: their periodic point divisors are distinct basis elements
of the divisor group of $\C/\Z$, so the translates together with all of their derivatives are
algebraically independent over the field generated by $z$ and the one-periodic meromorphic
functions.  In particular they are independent over $\C(z)$.

For one differential indeterminate $Y$, the successive derivatives of $Y'/Y$ have the triangular
form
\[
  \left(\frac{Y'}Y\right)^{(r)}
    =\frac{Y^{(r+1)}}Y+R_r\!\left(
      Y,\ldots,Y^{(r)}\right),
\]
where $R_r$ is rational.  Hence the change from the Gamma jets to $Y$ together with the jets of
$Y'/Y$ is birational and triangular.  A differential algebraic relation among a finite subfamily
of the digamma palette would therefore give an algebraic relation among the corresponding Gamma
jets, which is impossible.  Differential algebraic dependence is finitary, proving (i).

Gauss's expansion gives, locally normally off the pole set,
\begin{equation}\label{eq:gauss-digamma-series}
  G_\star(z)=\sum_{n=0}^{\infty}
    \left(\frac1{n+1}-\frac1{n+z+\tfrac12}\right).
\end{equation}
For $z=\sigma+it$ with $t>0$, each summand has imaginary part
\[
  \frac{t}{(n+\sigma+\tfrac12)^2+t^2}>0.
\]
Their sum is strictly positive, which proves (ii).

Equation~\eqref{eq:finite-digamma-orientation} differs from $G_\star$ by a rational function.
After any shift $\tau_q$ the difference still belongs to $\C(z)$.  Thus base translation preserves
both differential independence and the generated embedded field.  The pole at
$-\nu(p)-\tfrac12$ has residue $-1$ before reversal and $+1$ after reversal.  When $J$ is nonempty,
isolate a reversed pole.  All other terms are holomorphic and bounded in a sufficiently small
punctured neighbourhood, whereas vertical approach from the upper half-plane makes the imaginary
part of the reversed atom tend to $-\infty$.  Thus $G_{\star,J}$ is not Herglotz.  The case
$J=\varnothing$ is (ii), proving (iii).

On a compact set avoiding the poles, the summand in \eqref{eq:signed-digamma-completion} is
uniformly $O(n^{-2})$, because
\[
  \frac1{n+1}-\frac1{n+z+\tfrac12}
   =\frac{z-\tfrac12}{(n+1)(n+z+\tfrac12)}.
\]
The Weierstrass test proves local normal convergence for every sign sequence and also proves the
stated convergence of finite truncations.  Termwise isolation gives the residue
$-\varepsilon(q_n)$.  If all signs are positive, \eqref{eq:signed-digamma-completion} is
\eqref{eq:gauss-digamma-series} and is Herglotz.  If one sign is negative, its pole contributes
$-1/t$ on vertical approach while the regular remainder stays bounded, so the function is not
Herglotz.  This proves (iv).

Finally, (i) is precisely the source hypothesis used by the facet-sewing construction; (iii)
preserves it in every finite chamber.  The localized Stanley--Reisner realization therefore
applies verbatim, and (iv) supplies the asserted inverse-limit representative and pointing.
\end{proof}

\begin{theorem}[Canonical normalized source]\label{thm:canonical-digamma-source}
The function $G_\star$ is the unique meromorphic function $G$ on $\C$, holomorphic on the open
right half-plane, that satisfies
\begin{equation}\label{eq:canonical-digamma-cocycle}
  G(z+1)-G(z)=\frac1{z+\tfrac12}
  \qquad (\Re z>0)
\end{equation}
and the logarithmic normalization
\begin{equation}\label{eq:canonical-digamma-normalization}
  G(x)-\log x\longrightarrow\gamma
  \qquad (x\to+\infty).
\end{equation}
Thus the explicit Herglotz source is forced by a shift cocycle and a boundary normalization; it is
not an arbitrary representative chosen from a generic phase.
\end{theorem}

\begin{proof}
The digamma recurrence gives \eqref{eq:canonical-digamma-cocycle} for $G_\star$, while the classical
asymptotic $\psi(z)=\log z+O(z^{-1})$ in sectors away from the negative real axis gives
\eqref{eq:canonical-digamma-normalization} \cite[\S5.11]{NIST2010}.

Let $G$ be a second solution and put $H=G-G_\star$.  Equation
\eqref{eq:canonical-digamma-cocycle} makes $H$ one-periodic on the right half-plane.  For every
fixed $x>0$ and every $n\in\N$,
\[
  H(x)=H(x+n).
\]
Both normalized asymptotics show that the right-hand side tends to zero as $n\to\infty$; hence
$H(x)=0$ for every $x>0$.  The identity theorem gives $H=0$ on the right half-plane, and the
identity principle for meromorphic functions gives equality on $\C$.
\end{proof}

\begin{corollary}[Intrinsic source selection and infinite-field boundary]
\label{cor:resolved-boundaries}
No generic choice is required for the Herglotz product geometry: $G_\star$, the shifts $\tau_q$,
and every finite chamber are given by explicit formulas.  Theorem~\ref{thm:canonical-digamma-source}
makes this selection canonical, while $G_\star$ is also the unique Herglotz member of the signed
family $(G_\varepsilon)_\varepsilon$.  Every infinite orientation has its canonical analytic limit
$G_\varepsilon$.  Independently, its raw differential-field coordinate in the prime--Dirichlet
realization is exactly the projective tail class
\[
  [\varepsilon]\in
  \{\pm1\}^{\mathcal P}/(\text{finite changes and global sign}).
\]
Thus the canonical arithmetic source field $\mathscr L_{\boldsymbol{+}}$ contains exactly the
finite and cofinite reversal chambers, while every genuinely mixed infinite tail creates a
transverse differential-field direction.
\end{corollary}

\begin{proof}
The explicit Herglotz assertions are Theorems~\ref{thm:digamma-source} and
\ref{thm:canonical-digamma-source}.  The field classification is
Corollaries~\ref{cor:complete-palette-field-quotient} and
\ref{cor:finite-cofinite-source-field}.
\end{proof}

\begin{theorem}[Unique positive-definite pointing]\label{thm:positive-definite-pointing}
Fix $\sigma>1$.  The vertical trace
\[
  t\longmapsto P_\varepsilon(\sigma+it)
\]
is positive definite on $\R$ if and only if $\varepsilon=\boldsymbol{+}$.
\end{theorem}

\begin{proof}
It is the Fourier transform, up to the harmless choice of Fourier-sign convention, of the finite
signed measure
\[
  \mu_{\varepsilon,\sigma}
   =\sum_p\varepsilon(p)p^{-\sigma}\delta_{\log p}.
\]
For the all-positive orientation this is a positive measure, so its transform is positive
definite.  Conversely, Bochner's theorem \cite{Bochner1933} represents every continuous positive-definite function
as the Fourier transform of a finite positive measure.  Uniqueness of Fourier transforms of finite
measures identifies that measure with $\mu_{\varepsilon,\sigma}$.  Every atomic mass
$\varepsilon(p)p^{-\sigma}$ must therefore be positive, forcing every sign to be $+1$.
\end{proof}

\begin{corollary}[Closed-form constraint--orientation product]
\label{cor:closed-form-product}
Let $S\subset\mathcal P$ be finite and $J\subseteq S$, and let $\varepsilon_{S,J}$ be negative
exactly on $J$.  Starting from the single closed-form source $P_+$ and the fixed code
\eqref{eq:explicit-shift-code}, the complete palette of $P_{\varepsilon_{S,J}}$ is differentially
algebraically independent, generates $\mathscr L_+$, and realizes every finite localized
Stanley--Reisner differential matroid.  Its vertical trace is positive definite exactly when
$J=\varnothing$.  Hence the compatible positive vertices of all finite cubes converge to the
unique positive-definite point $\boldsymbol{+}$ of the full orientation cube.
\end{corollary}

\begin{proof}
Theorem~\ref{thm:explicit-source} gives independence and constraint realization.  The orientation
$\varepsilon_{S,J}$ differs from $\boldsymbol{+}$ on the finite set $J$, so
Corollary~\ref{cor:complete-palette-field-quotient} gives
$\mathscr L_{\varepsilon_{S,J}}=\mathscr L_+$.  Theorem~\ref{thm:positive-definite-pointing} gives
the last two assertions.
\end{proof}

\section{Canonical objects and scope}\label{sec:scope}

The three realizations expose the same product geometry at complementary levels.
\begin{enumerate}
  \item The residual Cauchy model proves that the Herglotz product occupies a robust generic phase,
    rather than being an isolated special-function coincidence.
  \item The digamma model supplies the closed-form Herglotz source $G_\star$, an explicit universal
    palette, and a canonical normally convergent representative of every infinite orientation.
  \item The prime--Dirichlet model supplies the exact differential-field stratification of the
    full orientation cube.
  \item Facet sewing records an arbitrarily prescribed finite hereditary constraint geometry in
    both explicit source palettes.
  \item The inverse-limit orientation cube carries two intrinsic positive pointings: Herglotz
    positivity in the Cauchy and digamma models, and vertical positive definiteness in the
    Dirichlet model.
  \item Its raw arithmetic field coordinate is the projective tail quotient
    \[
      \{\pm1\}^{\mathcal P}/(\text{finite changes and global sign}).
    \]
    Every fibre is countable and dense, and every nonconstant separated readout of the field class
    is nowhere continuous in the finite-chamber product topology.
  \item Its compact-open common-atom completion is constant across the entire source--orientation
    product, including every infinite orientation.
\end{enumerate}
Thus the infinite same-field boundary is no longer left open for the arithmetic fields generated
by the complete universal palettes.  Raw equality holds precisely along an eventual-equality or
eventual-opposition class; transverse tail directions contribute a jointly independent pair of
complete palettes.  Theorem~\ref{thm:projective-tail-field-bundle} packages these fields as the
canonical bundle over the projective-tail quotient.  In the Cauchy realization, compact-open
completion makes all source and orientation fibres coincide.
This is a positive structural law: the inverse limit is a projective-tail stratification of raw
differential fields over one completed analytic algebra, and its unique positive-definite section
meets one distinguished point of that stratification.

Theorem~\ref{thm:projective-tail-opacity} also identifies the exact cost of choosing raw-field
coordinates.  The orbit set has no nontrivial continuous separated geometry inherited from finite
chambers, and no Borel transversal exists.  The canonical object is therefore the projective-tail
groupoid together with its positive pointing, not a closed-form representative chosen from every
field fibre.  This is a necessity theorem rather than a defect of the construction.

In particular, Corollary~\ref{cor:finite-cofinite-source-field} settles the infinite-chamber
question without declaring all infinite reversals equivalent.  The canonical arithmetic source
field $\mathscr L_{\boldsymbol{+}}$ contains exactly the finite and cofinite reversal chambers.
Every mixed cofinal sign tail lies in a distinct field stratum and is separated from the canonical
arithmetic source by differential transcendence degree two at the raw-carrier level.

The pointing law applies directly to a proposed positivity argument whose input factors only
through an unpointed hypertranscendental field and finitely sewn differential constraints.  Such an
argument must add source-sensitive data capable of selecting the positive section.  In arithmetic
applications, the additional datum can be an Euler marking, a polarized kernel, or another rule
that distinguishes the all-positive tail class.  Theorem~\ref{thm:positive-definite-pointing}
provides exactly such a rule on the half-plane of absolute convergence.

The result does not assert the Riemann hypothesis.  The prime-zeta source is related to the Euler
product by
\[
  \log\zeta(s)=\sum_{k\geq1}\frac1kP_+(ks),\qquad \Re s>1,
\]
but positivity on this half-plane does not locate the singularities created by meromorphic
continuation.  A route to the Riemann hypothesis must therefore preserve the positive Euler
orientation while adding the functional equation and the continued divisor.  What the present
theorem rules out is more precise: no route invariant under finite Euler gauges can recover the
positive orientation from raw differential-field data alone, because it sees at most the
projective tail class.  Conversely, every viable route of this type must carry a marking that
survives the passage from finite chambers to the tail quotient and remains compatible with the
continued divisor.

The two earlier limitations are therefore replaced by exact positive statements.  First,
$G_\star$ is the unique cocycle-normalized source and the unique Herglotz member of its signed
family, so no generic source is selected.  Second, the canonical inverse-limit field bundle must
be projectively stratified before completion: its arithmetic raw field is constant exactly on the
eventual-equality and eventual-opposition class, while every genuinely mixed tail supplies a
transverse field direction.  Every stratum belongs to the same compact-open completed Cauchy
algebra.  The explicit models are complementary: the digamma model canonically realizes the
Herglotz geometry, the summable Cauchy model identifies the common completion, and the arithmetic
model reads the projective tail field coordinate.  Nothing is left as an unspecified same-field
boundary.

\begin{thebibliography}{99}

\bibitem{Bayart2018}
F. Bayart,
\emph{Hypercyclic algebras},
arXiv:1804.01730 (2018).

\bibitem{BesEtAl2017}
J. B\`es, J. A. Conejero, and D. Papathanasiou,
\emph{Hypercyclic algebras for convolution and composition operators},
arXiv:1706.08022 (2017).

\bibitem{Birkhoff1929}
G. D. Birkhoff,
\emph{D\'emonstration d'un th\'eor\`eme \'el\'ementaire sur les fonctions enti\`eres},
Comptes Rendus de l'Acad\'emie des Sciences de Paris \textbf{189} (1929), 473--475.

\bibitem{Bochner1933}
S. Bochner,
\emph{Monotone Funktionen, Stieltjessche Integrale und harmonische Analyse},
Mathematische Annalen \textbf{108} (1933), 378--410.

\bibitem{DoughertyJacksonKechris1994}
R. Dougherty, S. Jackson, and A. S. Kechris,
\emph{The structure of hyperfinite Borel equivalence relations},
Transactions of the American Mathematical Society \textbf{341} (1994), 193--225.

\bibitem{Hardouin2008}
C. Hardouin,
\emph{Hypertranscendance des syst\`emes aux diff\'erences diagonaux},
Compositio Mathematica \textbf{144} (2008), 565--581.

\bibitem{HassiWietsma2013}
S. Hassi and H. L. Wietsma,
\emph{Products of generalized Nevanlinna functions with symmetric rational functions},
arXiv:1312.7248 (2013).

\bibitem{Holder1886}
O. H\"older,
\emph{Ueber die Eigenschaft der Gammafunction keiner algebraischen Differentialgleichung zu
gen\"ugen},
Mathematische Annalen \textbf{28} (1886), 1--13.

\bibitem{Kolchin1973}
E. R. Kolchin,
\emph{Differential Algebra and Algebraic Groups},
Academic Press, New York, 1973.

\bibitem{Nedic2020}
M. Nedi\'c,
\emph{Characterizations of Herglotz--Nevanlinna functions using positive semi-definite functions
and the Nevanlinna kernel in several variables},
arXiv:2012.05532 (2020).

\bibitem{NIST2010}
F. W. J. Olver, D. W. Lozier, R. F. Boisvert, and C. W. Clark (eds.),
\emph{NIST Handbook of Mathematical Functions},
Cambridge University Press, Cambridge, 2010; online companion,
\emph{NIST Digital Library of Mathematical Functions}, \S5.11.

\bibitem{Pong2016}
W. Y. Pong,
\emph{Applications of differential algebra to algebraic independence of arithmetic functions},
Acta Arithmetica \textbf{172} (2016), 149--173.

\bibitem{RovnyakSakhnovich2004}
J. Rovnyak and L. A. Sakhnovich,
\emph{On the Krein--Langer integral representation of generalized Nevanlinna functions},
Electronic Journal of Linear Algebra \textbf{11} (2004), 1--15.

\bibitem{ShapiroSparer1986}
H. N. Shapiro and G. H. Sparer,
\emph{On algebraic independence of Dirichlet series},
Communications on Pure and Applied Mathematics \textbf{39} (1986), 695--745.

\bibitem{Stanley1975}
R. P. Stanley,
\emph{The upper bound conjecture and Cohen--Macaulay rings},
Studies in Applied Mathematics \textbf{54} (1975), 135--142.

\bibitem{Tsirivas2014}
N. Tsirivas,
\emph{Existence of common hypercyclic vectors for translation operators},
arXiv:1411.7815 (2014).

\end{thebibliography}

\end{document}
